% Pursuing Stacks (cotes 134-2 to 134-8): archive cote 134-5, site page 45 (PDF page 46), typescript page 232.
% typescript. Draft compared: ps4.tex:2202-2277.
% AI-transcribed by Claude Opus 5.5 (claude-opus-5-5, adaptive thinking, effort high), Message Batches API, batch ps2 (msgbatch_01NVwr1Yuo8qJf5rykeGBXPK), 2026-11-12; reviewed against both supplied scan halves; see CHECK_DECISIONS.tsv.
\scanpage{134-5/45}
\mspage{232}

defines an asphericity structure \xout{xx} $M_{as}$ on $M$, admitting $i$ as a ``\hand{\uline{\textcolor{black}{testing}}}%
\margin{X}
\hand{\uline{\textcolor{black}{functor}}}'' (i.e. a functor satisfying (i) and (ii) above), provided only
we assume that

(3) For any $a$ in $A$, $i^{*}(i(a))$ is an aspheric object of $A^{\wedge}$ .

This latter condition is automatically satisfied if $i$ is fully faithful,
for instance, if it is the inclusion functor of a full subcategory $A$
of $M$. Thus, \hand{\uline{\textcolor{black}{any small full subcategory}}} $A$ \hand{\uline{\textcolor{black}{of}}} $M$ \hand{\uline{\textcolor{black}{defines an asphericity}}}
\hand{\uline{\textcolor{black}{structure on}}} $M$, and we'll see in a minute that any asphericity structure
on $M$ can be defined this way.

\hand{\uline{\textcolor{black}{Proposition}}} 1 Let $M_{as}$ be an asphericity structure on $M$, and $A,B$ ob-%
jects in (Cat), and
\[ A \hand{\xrightarrow{\textcolor{black}{f}}} B \hand{\xrightarrow{\textcolor{black}{j}}} M_{\xout{xx}} \quad \xout{x} \]
be functors, with $j$ factoring through $M_{as}$ .

a) If $f$ is aspheric, then j is a testing functor iff $i=jf$ is .

b) If $j$ is fully faithful, and if \xout{bothxjxand} $i=jf$ \xout{are}\inserted[typed]{is a} testing func-%
tor\xout{\unsure{s}}, then $f$ is aspheric, and $j$ is a testing functor too.

Proof. a) follows trivially from
\[ B^{\wedge}_{as} = (f^{*})^{-1}(A^{\wedge}_{as}) , \]
which is one of the ways of expressing that $f$ is as\over{\unsure{s}}{p}heric\xout{\unsure{,}} (criterion
(ii') of prop.1, p. 214). And the first assertion in b) is a trivial
consequence of the definition of testing functors, and of aspher\xout{x}icity
of $f$ (by criterion (iii) on p. 214) -- and the second assertion of b)
now follows from a).

\hand{\uline{\textcolor{black}{Corollary}}} Let $i$: $A\hand{\longrightarrow} M$ be a testing functor for $(M,M_{as})$, and let
$B$ any \hand{\uline{\textcolor{black}{small}}} \inserted[typed]{full} subcategory of $M_{as}$ containing $i(A)$. Then the induced functor
$A\hand{\longrightarrow} B$ is aspheric, and the inclusion functor $B\hand{\hookrightarrow} M$ is a testing functor.

\hand{\uline{\textcolor{black}{R\xout{m}emarks}}} This shows, as announced above, that any asphericity structu-%
re on $M$ can be defined by a small full subcategory of $M$ (for instance,
the smallest full subcategory of $M$ containing $i(A)$). We have been slight-%
ly floppy though, while we defined testin\over{\unsure{g}}{g}\xout{xxx} functors \xout{xx} by insisting
that the source should be in (Cat) (which at times will be convenient),
whereas the $B$ we got here is merely small, namely \hand{\uline{\textcolor{black}{isomorphic}}} to an ob-%
ject of (Cat), but not necessarily in (Cat) itself. This visi\xout{x}bly is an
``unessential floppyness'', which could be straightened out trivially, ei-%
ther by enlarging accordingly the notion of testing functor, or by sub-%
mitting $M$ to the (somewhat artificial, admittedly) restriction that all

\margin{45}
